\documentclass{article}
\usepackage[UTF8]{ctex}
\usepackage{amsmath,amssymb,amsfonts}
\usepackage{enumitem}
\usepackage{geometry}
\geometry{a4paper, margin=1in}
% === 错题本排版优化（不参与题目标识） ===
\setlength{\parindent}{0pt}
\setlength{\parskip}{0.35em plus 0.10em minus 0.08em}
\setlength{\emergencystretch}{3em}
\setlist{itemsep=0.25em, topsep=0.40em, parsep=0pt}
\allowdisplaybreaks[2]
% === 错题本排版优化结束 ===
\begin{document}
\section*{河海大学2026年高等代数真题}
\section*{88．河海大学2026年高等代数试卷}
一．填空题．\\

\subsection{题目1}
1．设 \(A=\left(\begin{array}{lll}1 & 0 & 1 \\ 0 & 2 & 0 \\ 1 & 0 & 1\end{array}\right)\) ，若 \(n \geq 2\) 为正整数，则 \(A^{n}-2 A^{n-1}=\)

\subsection{题目2}
2．设 \(A, B\) 为三阶矩阵，\(|A|=3,|B|=2,\left|3 A^{-1}+2 B\right|=2\) ，则 \(\left|2 A+3 B^{-1}\right|=\) \(\_\_\_\_\)．

\subsection{题目3}
3．设三阶矩阵 \(A=\left(\begin{array}{lll}1 & 2 & -2 \\ 2 & 1 & 2 \\ 3 & 0 & 4\end{array}\right)\) ，三维列向量 \(\alpha=(t, 1,1)^{\mathrm{T}}\) ，已知 \(A \alpha\) 与 \(\alpha\) 线性相关，则 \(t=\)\\
\(\_\_\_\_\)．

\subsection{题目4}
4．设 \(n\) 维向量 \(\alpha=(t, 0, \cdots, 0, t)^{\mathrm{T}}, t \neq 0, E\) 是 \(n\) 阶单位矩阵，矩阵 \(A=E-\alpha \alpha^{\mathrm{T}}, B=E+\frac{1}{t} \alpha \alpha^{\mathrm{T}}\) ，其中 \(A\) 的逆矩阵为 \(B\) ，则 \(t=\) \(\_\_\_\_\)．

\subsection{题目5}
5．设 \(A\) 是复数域 C 上的 \(n\) 阶方阵，特征多项式为 \(|\lambda E-A|=(\lambda-a)^{n-1}(\lambda-b), a, b\) 是两不等的复数．若 \(A\) 的任意三个特征向量都是线性相关的，则 \(A\) 的若尔当标准形为 \(\_\_\_\_\)．
\end{document}
